\section{Discussion and conclusion}

ZoomST makes observation granularity an explicit variable in spatial transcriptomic pretraining by jointly learning granularity-specific tissue representations and transformations between them. Its design combines molecular--spatial dual-view alignment within each granularity with soft distributional equivariance across nested neighborhoods. A shared neural ODE learns continuous transformations from supervision at discrete granularities, while the representations at each granularity remain separately usable. Different neighborhood extents favor different tissue structures; across the held-out MOSTA sections, finer contexts yield higher overall anatomical agreement, while coarser contexts retain more physical neighbors. The scale-specific comparison with the matched-update No Cross ablation shows that learning cross-granularity transformations enriches developmental-stage ordering in coarse joint representations while retaining comparable tissue-identification performance. Learned relations between observation scales thus help broad-context embeddings retain information about developmental progression alongside tissue identity and local spatial organization.

The learned transformations also predict directly encoded coarse and unseen intermediate states. Across held-out MOSTA sections, predicted responses generally track relative differences in encoded change, with anatomical analyses linking these differences to neighborhood composition. This makes observation granularity both a source of self-supervision and an axis for tissue analysis.
